% !PRINTED 624
\weberpage{624}
% Running head "624  Vierundzwanzigster Abschnitt.  § 171." is page furniture and
% is not transcribed as body text; only the printed number is recorded above.
%<<CONT
% (This page's own footnote apparatus begins mid-sentence, continuing the
% footnote 1) that was left open at the bottom of the previous page: "...Clebsch,
% Über die" + "Anwendung der Abelschen Funktionen..." joins into one sentence.
% The main-text continuation below is a fresh section, not a continued sentence.)

\begin{center}
\textbf{§ 171. Definition der algebraischen Funktionen.}
\end{center}

Eine Variable $\theta$ heißt eine algebraische Funktion einer
unabhängigen Veränderlichen $z$, wenn die beiden Variablen durch eine
algebraische Gleichung
\weq{1}{F(\theta,z)=0}
miteinander verbunden sind. $F$ bedeutet hierin einen Ausdruck von der
Form
\weq{2}{F(\theta,z)=a_0\theta^n+a_1\theta^{n-1}+\cdots+a_{n-1}\theta+a_n,}
dessen Koeffizienten $a_0, a_1, \ldots, a_n$ ganze rationale Funktionen
von $z$ ohne gemeinschaftlichen Teiler sind. Über die konstanten
Koeffizienten in diesen Ausdrücken machen wir weiter keine
Voraussetzung, als daß es reelle oder komplexe Zahlen
sind\textsuperscript{1)}.

Wir setzen dabei die Funktion $F(\theta,z)$
in dem Sinne als irreduzibel voraus, daß sie nicht in Faktoren
zerfallen soll, die selbst rationale Funktionen von $\theta$ und $z$
sind.

Jede ganze Funktion $G(\theta,z)$ läßt sich nach Bd.\,I, § 20
nur auf eine Weise in irreduzible Faktoren zerlegen (abgesehen von
konstanten Faktoren). Wir sagen, daß die Funktion $G(\theta,z)$
%CONT>>
% body text is cut off here by the page break; the sentence continues on the
% next page, which is not included in this scan set.

\smallskip
\noindent\rule{5em}{0.4pt}

\noindent Anwendung der \gesperrt{Abelschen} Funktionen in der Geometrie, Crelles
Journ., Bd.\,63, 1864. \gesperrt{Clebsch} und \gesperrt{Gordan},
Theorie der \gesperrt{Abelschen} Funktionen, Leipzig 1866. \gesperrt{Brill} und
\gesperrt{Noether}, Über die algebraischen Funktionen und ihre
Anwendung in der Geometrie, Mathematische Annalen VII, 1874.
\gesperrt{Brill} und \gesperrt{Noether}, Die Entwickelung der Theorie
der algebraischen Funktionen, Bericht der Deutschen
Mathematiker-Vereinigung 1894. \gesperrt{Appell} et \gesperrt{Goursat},
Théorie des fonctions algébriques, Paris 1895. Über die
zahlentheoretischen Methoden: \gesperrt{Kronecker}, Über die
Diskriminante algebraischer Funktionen, Crelles Journ., Bd.\,91, 1881,
und: Festschrift zu \gesperrt{Kummers} Doktor-Jubiläum, Crelles
Journ., Bd.\,92, 1881. \gesperrt{Dedekind} und \gesperrt{Weber},
Theorie der algebraischen Funktionen einer Veränderlichen, Crelles
Journ., Bd.\,92, 1879. \gesperrt{Hensel} und \gesperrt{Landsberg},
Theorie der algebraischen Funktionen einer Variablen, Leipzig 1902.
Die Form der Theorie, die im folgenden dargestellt ist, stützt sich
auf den in Bd.\,II, § 153 eingeführten Begriff der Funktionale, die in
dieser Form bereits in einer von \gesperrt{Wellstein} angeregten
Straßburger Dissertation von \gesperrt{Rehfeld} (1906) zur Behandlung
eines speziellen Problems der Funktionentheorie angewandt wurden.

\smallskip
\noindent\textsuperscript{1)}~Es würde nirgends eine Lücke bleiben,
wenn wir uns dabei auf \gesperrt{algebraische Zahlen} beschränken wollten.
Betrachtet man $\theta, z$ als Cartesische Koordinaten in der Ebene,
so stellt die Gleichung (1) eine algebraische Kurve dar, und man kann
die Theorie dieser Kurven anwenden, wie es \gesperrt{Clebsch} und
\gesperrt{Gordan} und \gesperrt{Brill} und \gesperrt{Noether} getan
haben. Freilich haben dabei nur die reellen Werte dieser Variablen
eine wirklich anschauliche Bedeutung.
